\documentclass{article}

\usepackage{../../apxproof}
\usepackage{zref-clever}

% Make the hook visible, so that its argument is part of what is compared.
\renewcommand{\apxproofhook}[1]{[hook: \texttt{#1}]\ }

\newtheoremrep{theorem}{Theorem}
\newtheorem{lemma}{Lemma}

\begin{document}

\section{Labels set with the standard command}

Here labels are set with \verb|\label|, which \textsf{zref-clever} picks up
through the kernel's label hook.

\begin{theoremrep}
  \label{th:hook}%
  A theorem, with an equation:
  \begin{equation}
    \label{eq:hook}
    1=1
  \end{equation}
\end{theoremrep}

\begin{proof}
  Deferred proof, referring to \zcref{eq:hook}.
\end{proof}

References: \zcref{th:hook}, \zcref{eq:hook}, and the appendix copy of the
equation, \zcref{eq:hook-apx}.

\section{Labels set with the zref command}

\zcsetup{labelhook=false}

\begin{theoremrep}
  \zlabel{th:z}%
  Another theorem, with an equation:
  \begin{equation}
    \zlabel{eq:z}
    2=2
  \end{equation}
\end{theoremrep}

\begin{proof}
  Deferred proof, referring to \zcref{eq:z}.
\end{proof}

References: \zcref{th:z}, \zcref{eq:z}, and the appendix copy of the
equation, \zcref{eq:z-apx}.

\section{Proof titles naming their result}

\zcsetup{labelhook=true}

\begin{lemma}
  \label{lem:one}%
  A lemma whose proof is not deferred.
\end{lemma}

\begin{lemma}
  \label{lem:two}%
  Another one.
\end{lemma}

\begin{proof}[Proof of \zcref{lem:one}]
  Inline proof, whose title names its result.
\end{proof}

\begin{proof}[Proof of \zcref*{lem:two}]
  Inline proof, whose title names its result in the starred form.
\end{proof}

\end{document}
